% ============================================================
\section{Data Sources}
\label{sec:data}
% ============================================================

All datasets used by the pipeline are freely accessible from global
open repositories.
Table~\ref{tab:datasources} lists each source with the relevant
variables, spatial resolution, and access method.
No registration is required for any source; all downloads are executed
automatically at run time through standard HTTP requests.

\begin{table}[H]
\centering
\caption{Data sources used by the CRNS site characterization pipeline.}
\label{tab:datasources}
\small
\begin{tabularx}{\linewidth}{lXll}
\toprule
Source & Variables & Resolution & Access \\
\midrule
Copernicus DEM GLO-30
  & Elevation (m a.s.l.)
  & 30 m
  & AWS S3 (open) \\
ESA WorldCover 2021 v200
  & Land-cover class (11 classes)
  & 10 m
  & AWS S3 / Planetary Computer \\
OpenStreetMap (Overpass API)
  & Water bodies, roads, buildings
  & Vector
  & REST API \\
SoilGrids 2.0
  & Sand, silt, clay (\%), bulk density, SOC
  & 250 m
  & ISRIC REST API \\
ERA5-Land / Open-Meteo Archive
  & Temperature, RH, snowfall, snow depth (daily, 5 yr);
    soil moisture 0--100 cm (hourly, $\geq$2015)
  & $\sim$9 km (soil moisture) / $\sim$30 km (climate)
  & Open-Meteo REST API \\
MODIS MOD15A2H
  & Leaf Area Index (LAI), 8-day composites
  & 500 m
  & NASA AppEEARS / pystac \\
Macrostrat API v2
  & Geological units, lithology, age
  & Global (variable)
  & Macrostrat REST API \\
\citet{smart2019} / \texttt{crnpy}
  & Geomagnetic rigidity cutoff $R_c$
  & $1^\circ \times 1^\circ$ grid
  & Python package \\
\bottomrule
\end{tabularx}
\end{table}

\subsection{Digital Elevation Model}

The Copernicus DEM GLO-30 \citep{copernicus_dem} provides a
near-global 30-m resolution surface elevation model derived from
TanDEM-X radar interferometry.
The pipeline downloads the tiles covering a square bounding box
$[{-}r_\mathrm{buf},\,{+}r_\mathrm{buf}]$ centred on the site,
with default buffer $r_\mathrm{buf} = 3\,\rzz \approx \SI{350}{m}$.
Tiles are mosaicked and reprojected to a metric coordinate system
(UTM, WGS-84) before use.

The DEM is the master grid: all raster products are resampled to
match its spatial extent and pixel spacing.

\subsection{Land-Use / Land-Cover}

ESA WorldCover 2021 \citep{worldcover2021} provides a 10-m global
land-cover classification based on Sentinel-1 and Sentinel-2
imagery.
The 11 classes (tree cover, shrubland, grassland, cropland,
built-up, bare/sparse vegetation, snow/ice, permanent water,
herbaceous wetland, mangroves, moss/lichen) are aggregated to the
hydrogen equivalence groups of Table~\ref{tab:fH}.

OpenStreetMap (OSM) features — surface water bodies, roads with
sealed surfaces, and building footprints — are queried via the
Overpass API within the $2\,\rzz$ radius of the footprint.
OSM vectors complement WorldCover for fine-scale features
(individual buildings, canals) that may be below the 10-m raster
resolution.

The JRC Global Surface Water dataset \citep{pekel2016} provides
occurrence frequency (0–100\%) for permanent and seasonal surface
water, used to distinguish permanent from temporary water bodies
in the $f_H$ assignment.

\subsection{Soil Properties}

SoilGrids 2.0 \citep{soilgrids2020} provides global predictions at
six depth intervals (0–5, 5–15, 15–30, 30–60, 60–100, 100–200 cm)
for:
\begin{itemize}[nosep]
  \item sand, silt, clay fractions [\si{\gram\per\kilo\gram}];
  \item bulk density (fine earth) [\si{\centi\gram\per\cubic\centi\metre}];
  \item soil organic carbon (SOC) content [\si{\deci\gram\per\kilo\gram}].
\end{itemize}
The pipeline uses the $0$–$30\,\si{cm}$ depth-weighted mean, which
matches the typical $z_{86}$ penetration range.
Lattice water $\lw$ is derived from the clay fraction using the
empirical relation of \citet{kohli2021}:
\begin{equation}
  \lw = 0.097\,f_\mathrm{clay} + 0.033,
  \label{eq:lw_clay}
\end{equation}
where $f_\mathrm{clay}$ is the clay mass fraction [\si{g/g}].

\subsection{Atmospheric Data}

Monthly climatologies of near-surface air temperature, relative
humidity, snowfall, and snow depth are downloaded from the
Open-Meteo ERA5 archive \citep{hersbach2020} for the 5 years
preceding the analysis date.
Daily values are aggregated to monthly means and then averaged
across years to obtain a 12-month climatology.

From temperature $T$ [\si{\celsius}] and relative humidity
$\mathrm{RH}$ [\%], the atmospheric water-vapour density
$\rho_{WV}$ [\si{\gram\per\cubic\metre}] is computed via the
Magnus formula:
\begin{equation}
  e_\mathrm{sat}(T) = 611.2 \exp\!\left(\frac{17.67\,T}{T + 243.5}\right)
  \quad[\si{\pascal}],
  \label{eq:esat}
\end{equation}
\begin{equation}
  \rho_{WV} = \frac{\mathrm{RH}}{100}\,
              \frac{e_\mathrm{sat}}{R_v\,T_K}
  \quad[\si{\gram\per\cubic\metre}],
  \label{eq:rhoWV}
\end{equation}
where $R_v = \SI{461.5}{\joule\per\kilo\gram\per\kelvin}$ is the
specific gas constant for water vapour and $T_K$ is temperature
in Kelvin.

Snow-water equivalent (SWE) is estimated as
\begin{equation}
  \mathrm{SWE} = h_\mathrm{snow} \times \rho_\mathrm{snow},
  \label{eq:swe}
\end{equation}
with $h_\mathrm{snow}$ the mean snow depth [m] and
$\rho_\mathrm{snow} = \SI{250}{\kilo\gram\per\cubic\metre}$ a
representative seasonal density for wet alpine snow.

\subsection{Soil Moisture — ERA5-Land}
\label{sec:data:era5sm}

Hourly volumetric soil moisture is retrieved from the ERA5-Land
reanalysis via the Open-Meteo Historical API.
The module \texttt{era5sm.py} downloads three depth layers:
0--7\,cm, 7--28\,cm, and 28--100\,cm, all in
\si{\cubic\metre\per\cubic\metre}.
Downloads proceed year by year from a configurable
\texttt{start\_year} (default 2015) to the current date, with
each annual time-series stored as a compressed \texttt{.npz} file
(\texttt{hourly\_YYYY.npz}).
Monthly statistics (mean, standard deviation, minimum, maximum, and
observation count) are derived from the full multi-year record and
cached in a single \texttt{monthly\_agg.npz}.
On subsequent runs only missing years are downloaded; the monthly
aggregates are recomputed only if the set of cached years has
changed.

The three ERA5-Land depth layers bracket the typical CRNS
penetration depth $z_{86}$, making the 0--7\,cm layer the primary
reference for the data-fusion step (Section~\ref{sec:methods:smfusion}).

\subsection{Vegetation (LAI)}

The MODIS MOD15A2H product provides 8-day, 500-m composite
estimates of leaf area index (LAI) [\si{\metre\squared\per\metre\squared}].
Monthly means are derived from all cloud-free composites within
the analysis year.
LAI is used to estimate above-ground biomass hydrogen (AGBH)
as described in Section~\ref{sec:methods:agbh}.

\subsection{Geology}

The Macrostrat API v2 provides global geological map data at
multiple resolution scales (``large'', ``medium'', ``small'').
For each query point the pipeline retrieves the dominant
lithological unit(s), their Macrostrat rock type, geologic age,
and estimated hydrogen inventory impact on $\lw$.
Where lithological data are unavailable (ocean, ice sheet), the
pipeline returns a ``not available'' flag and uses the
SoilGrids-derived $\lw$ without geological supplement.
